% chapter I/introduction.tex

\chapter{Introduction}
\label{chap:introduction}

\section{Overview}
\label{sec:overview}

The rapid urbanization of developing countries has placed significant pressure on existing transportation systems, as growing populations rely heavily on automobiles for mobility, resulting in traffic volumes that exceed road capacity. When traffic demand exceeds the system's capacity, the transportation network becomes overloaded. Regional planning studies, specifically in the Calabarzon Traffic Management Plan 2017-2022,  identifies road sections that exceed road capacity as a major concern, noting that such conditions generate travel delays and chokepoints at intersections (Regional Development Council Region IV-A, 2017). This imbalance between increasing traffic flows and limited road capacity is commonly defined as congestion (Akthar \& Modipour, 2021). 

Traffic congestion is a growing problem in cities across the world and has a major impact on urban transport systems, particularly on highly populated areas, as it leads to considerable delays and costs (Faheem et al., 2024). Expanding road infrastructure is a common policy response to urban traffic congestion and has been widely implemented in countries such as the United States (Hansen et al., 1993), China, and India (Pucher et al., 2007). While it may seem that road expansions dilute existing traffic, it also induces additional travel demands which is likely to negate the expected benefits of the increased road capacity. When new roads are constructed, they induce additional traffic by facilitating car travel and incentivizing individuals to choose the optimal route (Chen \& Klaiber, 2020).

This behavior aligns with the Braess Paradox (BP), a critical theoretical framework for understanding how the addition of new road links can, under certain conditions, fail to improve traffic flow. First formulated by Dietrich Braess in 1968, the paradox postulates that adding capacity to a road network can, counter-intuitively increase the overall travel time when each individual in the network acts selfishly to minimize their own travel time (Rapoport et al., 2009). 

A challenge in infrastructure planning is whether an increased capacity always leads to an improved system performance. In transportation networks, the Braess Paradox (BP) demonstrates that under a selfishly routed system, increasing capacity can lead to an increased to higher equilibrium travel times (Mou \& Qin, 2025). Because of these complexities, planners utilize transportation models which are mathematical and computational representations of real-world transportation systems that are developed to simulate the interactions between network capacity and travel demand (Ortúzar \& Willumsen, 2024). These models are indispensable for forecasting future travel behaviors and demand, allowing city planners to evaluate the potential impact of infrastructure projects such as the closure or constructions of road sections (Barceló, 2010). 

Historically, transportation modeling in the Philippines has been predicated on the conventional Four-Step-Model (FSM), a forecasting framework utilized in urban transportation planning to predict network flows by aggregating travel data. Widely utilized in foundational studies such as the Metro Manila Urban Transportation Integration Study (MMUTIS), the FSM is a sequential process that comprises trip generation, trip distribution, mode split, and traffic assignment, treating traffic as a uniform, continuous flow rather than a collection of individual actors (Japan International Cooperation Agency, 1999).  Moreover, similar approaches have been used in the context of Iloilo City, specifically in the old Calle Real de Iloilo City, where circulation and public transport proposals have employed FSM frameworks in modeling the urban road network and evaluate alternative schemes (Alarcon et al., 2012). While the FSM provides an essential strategic overview for long term transportation planning, its aggregate nature struggles to capture the complexities of the Philippine Transport Systems. In particular, traditional models, like the FSM, often struggle to represent the unpredictability of paratransit modes like jeepneys and tricycles, whose flexible routes and irregular stopping behaviors do not align with fixed network definitions (Regidor, 2012). To address these limitations, urban planners increasingly use computational modeling to traffic flow and guide infrastructure development (Barcelo, 2010). 

Among modern methods, topological analysis utilizing computational frameworks has emerged as an effective tool. Specifically, the framework established by Youn et al. (2008) demonstrates how applying a mathematical cost function, like the Bureau of Public Roads (BPR) function, to empirical network data can be used to identify inefficiencies in real-world cities. To quantify this inefficiency, Youn et al. (2008) utilized the “Price of Anarchy” (PoA), a metric defined as the ratio between the total travel cost of a Nash Equilibrium and the total cost of a system-optimized social optimum. To identify the BP, the authors computed for the PoA to compare Nash equilibrium costs on the original network against the costs on the modified networks where individual streets were closed to traffic. By doing so, they were able to determine specific links in cities where the removal of a specific street counterintuitively decreased the overall travel delay in the Nash equilibrium, which would identify that street as a cause of BP. Consequently, this makes computational modeling the appropriate tool in investigating the Braess Paradox. By analyzing this discrepancy, the researchers can mathematically test whether the closure of specific road segments could counterintuitively improve overall traffic flow in a local setting. 
    
Iloilo City presents a compelling setting for applying this integrated perspective due to its status as the second highly urbanized city (HUC) in the Philippines (Bureau of Local Government Finance, 2025). Unlike the saturated network of Metro Manila, Iloilo City is currently at a critical urbanization inflection point, characterized by a surge of vehicular volume that exceeds its historic road network (Iloilo City Planning and Development Office, 2021; Ramos, 2025). In particular, the local government has pursued capacity expansion to accommodate increasing traffic demands, most notably through the construction of the Ungka Flyover. When traffic was re-routed to the adjacent at-grade intersections during the flyover's temporary closure, these nearby corridors absorbed higher volumes than they were designed for, leading to noticeable increases in congestion along those routes (Macalalag, 2023).

Similarly, Marzan and Castor (2025) report that a recent survey conducted by Jarobilla identified 38 traffic chokepoints across the city, many of which are found in primary roads such as Diversion Road, Jaro Plaza, and the City Proper. In the report, Jarobilla characterized chokepoints as road segments where congestion is obvious, including areas which require additional adjustments to improve the flow of traffic. These chokepoints emerge from geometric constraints, pedestrian activity, and roadside operations that constrain road capacity and intensify congestion. Local authorities note that the list of checkpoints requires continual updating as traffic conditions shift in response to changes in signal timing, operational adjustments, and the implementation of temporary road closures (Conlu, as cited in Philippine News Agency, 2022). Rather than simply reducing capacity, road closures often reconfigure the available network by redirecting vehicles to alternative road sections, effectively creating new or temporary routes that must absorb the redistributed traffic volumes from the reduced vehicles (Kiec et al., 2020).  

Despite the presence of chokepoints and the occurrence of temporary closures across the city, limited research has examined how road closures affect traffic flow in Iloilo City. Existing local studies predominantly focus on identifying congestion locations or analyzing intersection-specific volumes (Pragados, 2001) but do not explore how closures propagate through the network. While some computational studies have been conducted in the city, they have focused on mode choice behavior (Sosuan \& Fillone, 2014) or the feasibility of new transit systems like the Bus Rapid Transit (Calderon et al, 2014), rather than the dynamics of a driver's selfish route choice under a constrained network.  To address this gap, this research develops a computational network model that is designed to identify the presence of the Braess Paradox within Iloilo City's network. This topological analysis constructs a representation of the urban network across the city's major districts, specifically ity Proper, Jaro, La Paz, Lapuz, Mandurriao, Molo, and Villa-Arevalo. By evaluating traffic equilibrium at this scale, the research aims to establish a foundation for assessing how the effects of road closures propagate throughout a local network.

\section{Scope and Limitation}
\label{sec:scope_limitation}

The study evaluates the existing road network of Iloilo City and its susceptibility to congestion through the application of the Braess Paradox (BP). The network model was created based on the DPWH Atlas 2024, which focuses exclusively on the primary arterial road networks within Iloilo City. The scope of the network is also limited within Iloilo City's primary districts: City Proper, Jaro, La Paz, Lapuz, Mandurriao, Molo, and Villa-Arevalo.

The Braess Paradox was observed by developing a computational network model that utilizes topological analysis to evaluate traffic equilibrium. The study modeled route choices mathematically by comparing user-optimized traffic assignments against system-optimized configurations. The application of the BP was tested through the computational, iterative removal of road segments to determine if network efficiency counterintuitively improved.

Following the topological modeling framework of Youn et al. (2008), the development of the computational model relied on constructing a directed graph of the network topology using empirical road data. Specifically, Google Maps was utilized in pinpointing key intersections, which served as the network's exact mathematical nodes. The specific roads connecting these intersections, which are the directed edges or the links between local origin-destination points, were then defined by using structural data such as physical distance in meters, number of lanes, and directionality. By linking these nodes and edges, the study built a comprehensive representation of the city's transportation network. Together, these variables served as parameters for the Bureau of Public Roads (BPR) function.

Moreover, the model required the generation of predefined OD pairs to distribute traffic across the graph. The experiment computationally scaled the total traffic volume across these OD pairs. This scaling was necessary in order to maximize the Price of Anarchy (PoA) and identify the specific traffic volume where the network reached peak inefficiency. The iterative removal of roads was then performed under these maximized volume conditions to test the Braess Paradox at its theoretical limit.

 However, the transition of a complex road network into a discrete computational graph resulted in some research limitations. First, the model utilizes OD pairs, which were traffic demand projections for the year 2022, derived from the JICA (2004) socio-economic framework. Due to the lack of comprehensive traffic data of Iloilo City in recent years, the predicted OD pairs by JICA (2004) are one of the latest comprehensive traffic data of Iloilo City with high-resolution OD pairs available. Still, it establishes a temporal gap between the 2004 baseline and the current 2026 conditions that may be prone to bias or propagation of error in the modern trip distribution patterns.

 Second, the model excludes secondary, tertiary, and local roads in the creation of the network model. This limitation is due to the insufficient traffic data, such as the OD pair matrix, among the other roads. The shift of flow to the primary arteries may over-represent congestion levels on major links.

 Lastly, another key limitation is the exclusion of highly irregular driving behavior of jeepneys, particularly lane disruption caused by their “para” system, pedestrian crossings, traffic incidents, and weather conditions. By treating traffic as a continuous mathematical flow subject to the BPR function, the model may underrepresent highly localized, unpredictable congestion events that a physical simulation might capture.

\section{Problem Statement}
\label{sec:problem_statement}

Traffic congestion poses a significant barrier to a country’s economic growth, leading to environmental pollution and a decline in the quality of life (Fattah et al., 2022). In the Philippines, rapidly growing urban centers, including Iloilo City, have driven increased demand for transportation. According to Landoy et al. (2012, as cited in Sosuan \& Fillone, 2014), traffic congestion is already eminent along Iloilo City’s primary roads. A common intuitive solution to congestion is to expand capacity by building new roads or adding links. However, empirical evidence from decades of traffic data across the United States shows that increasing road capacity does not necessarily alleviate congestion (Stromberg, 2014). 

This global trend demonstrates that similar outcomes may also occur in developing urban centers such as Iloilo City. This counterintuitive phenomenon, known as the Braess Paradox, describes a situation in which strategically adding a new road link results in an overall increase in travel costs for all network users, rather than reducing them (Rapoport et al., 2009). The potential for such an outcome in Iloilo City highlights a critical gap between intuitive infrastructure planning and actual traffic network behavior. Thus, there is a pressing need to identify such paradoxical outcomes in a local setting. This study will address this need by developing and testing a model within a focused study area, which will serve as the foundational case study for the analysis. 


\section{Research Objectives}
\label{sec:research_objectives}

\subsection{General Objective}
\label{sec:general_objective}

The study aims to develop, test, and identify the existence of the Braess Paradox within Iloilo City's network, specifically the primary roads of the Iloilo City's seven districts: Villa-Arevalo, City Proper, Jaro, La Paz, Lapuz, Mandurriao, Molo. The results will provide a foundation for future city-wide analysis and contribute to more data-driven decisions in road planning and congestion management. 

\subsection{Specific Objectives}
\label{sec:specific_objectives}

Specifically, the study aims to:

\begin{itemize}
    \item Collect and process road network data, including road length, travel time, traffic volume, and other factors in the primary roads of the Iloilo City's seven districts: Villa-Arevalo, City Proper, Jaro, La Paz, Lapuz, Mandurriao, Molo.
    \item Construct a computational network model of Iloilo City using empirical structure and traffic data as parameters.
    \item Simulate various road closures to observe changes in overall travel time and traffic distribution.
    \item Identify and analyze routes affected by the Braess Paradox, in which adding or removing a road segment can counterintuitively improve network efficiency.
    \item Evaluate the effectiveness of road closure strategies in reducing travel time and improving traffic flow in the study area.
\end{itemize}

\section{Significance of the Study}
\label{sec:significance_of_the_study}

Addressing complex traffic networks is of paramount significance as it is crucial for growing urbanized areas such as Iloilo City (BLGF, 2025). Traffic congestion is already established along some of the primary roads in Iloilo City (Landoy, 2012, as cited in Sosuan \& Fillone, 2014). This congestion serves as more than an inconvenience; as stated by the (Metropolitan Transportation Commission, 2022), traffic congestion increases travel time and leads to more unpredictable estimates of arrival. Currently, there are few to no advanced traffic management systems integrated into the development and planning of the traffic networks in the city (Conlu, 2022). Thus, the application of the Braess Paradox to these networks provides a novel way to tackle congestion by analyzing scenarios where adding road capacity might actually cause congestion to persist or worsen (Braess, 2005). 

The significance of this study extends to improving the overall quality of life in the city and ensuring sustained economic growth (Fattah, 2022). For the planning and development offices of the provincial and city government, the simulation model serves as a low-cost option for testing modifications to the road network before proceeding with final physical implementations (ICPDO, 2021). For the general public and commuters, effective road closure strategies can lead to more predictable and shorter travel times, reducing the daily stress associated with traffic-congested commutes (Marve, 2016).

For the local economy, an optimized transportation network helps businesses lower logistics expenses, such as delivery times, and establishes a system that can withstand the growing urbanized economic demands of the city for future investment (Onyeneke, 2018). Additionally, a more efficient traffic network leads to reduced vehicle emissions, directly contributing to the environmental sustainability and efficiency of the city  (Marve, 2016).